\documentclass[11pt,a4paper]{article}

\usepackage[T1]{fontenc}
\usepackage[utf8]{inputenc}
\usepackage[italian]{babel}
\usepackage{lmodern}
\usepackage{microtype}
\usepackage[a4paper,margin=2.15cm,headheight=15pt]{geometry}
\usepackage{booktabs,longtable,tabularx,array}
\usepackage{ragged2e}
\usepackage{enumitem}
\usepackage{xcolor}
\usepackage[most]{tcolorbox}
\usepackage{fancyhdr}
\usepackage{hyperref}
\usepackage{xurl}
\usepackage{listings}
\usepackage{titlesec}
\usepackage{amsmath,amssymb}

\definecolor{ddtadark}{RGB}{28,38,52}
\definecolor{ddtagray}{RGB}{244,246,248}
\definecolor{ddtablue}{RGB}{42,88,135}
\definecolor{ddtagreen}{RGB}{48,111,78}
\definecolor{ddtaorange}{RGB}{148,91,25}
\definecolor{ddtared}{RGB}{140,48,48}
\definecolor{ddtapurple}{RGB}{92,63,122}

\hypersetup{
  colorlinks=true,
  linkcolor=ddtablue,
  urlcolor=ddtablue,
  citecolor=ddtablue,
  pdftitle={Documentation and Base Analysis Authoring Guide},
  pdfauthor={Documentation-Driven Threat Analysis research workspace}
}

\pagestyle{fancy}
\fancyhf{}
\lhead{DDTA -- Authoring Guide}
\rhead{R3}
\cfoot{\thepage}

\titleformat{\section}{\Large\bfseries\color{ddtadark}}{\thesection}{0.7em}{}
\titleformat{\subsection}{\large\bfseries\color{ddtadark}}{\thesubsection}{0.7em}{}
\titleformat{\subsubsection}{\normalsize\bfseries\color{ddtadark}}{\thesubsubsection}{0.7em}{}

\setlist[itemize]{leftmargin=1.55em,itemsep=0.22em,topsep=0.35em}
\setlist[enumerate]{leftmargin=1.75em,itemsep=0.22em,topsep=0.35em}
\setlength{\parskip}{0.35em}
\setlength{\parindent}{0pt}
\setlength{\emergencystretch}{3em}

\lstset{
  basicstyle=\ttfamily\small,
  columns=fullflexible,
  keepspaces=true,
  frame=single,
  framerule=0.3pt,
  rulecolor=\color{ddtablue},
  backgroundcolor=\color{ddtagray},
  xleftmargin=0.6em,
  xrightmargin=0.6em,
  aboveskip=0.5em,
  belowskip=0.5em,
  breaklines=true,
  showstringspaces=false
}

\newtcolorbox{closedbox}[1][REGOLA STABILE]{enhanced,breakable,colback=ddtagray,colframe=ddtagreen,title={#1},fonttitle=\bfseries}
\newtcolorbox{workingbox}[1][REGOLA BASELINE-SCOPED]{enhanced,breakable,colback=white,colframe=ddtaorange,title={#1},fonttitle=\bfseries}
\newtcolorbox{openbox}[1][NOT SPECIFIED / DIAGNOSTIC]{enhanced,breakable,colback=white,colframe=ddtapurple,title={#1},fonttitle=\bfseries}
\newtcolorbox{goodbox}[1][PREFERIBILE]{enhanced,breakable,colback=white,colframe=ddtagreen,title={#1},fonttitle=\bfseries}
\newtcolorbox{badbox}[1][DA EVITARE]{enhanced,breakable,colback=white,colframe=ddtared,title={#1},fonttitle=\bfseries}
\newtcolorbox{examplebox}[1][ESEMPIO FACIAL ACCESS]{enhanced,breakable,colback=ddtagray,colframe=ddtablue,title={#1},fonttitle=\bfseries}
\newtcolorbox{conceptbox}[1][CONCETTO]{enhanced,breakable,colback=white,colframe=ddtablue,title={#1},fonttitle=\bfseries}

\newcolumntype{Y}{>{\RaggedRight\arraybackslash}X}
\newcolumntype{L}[1]{>{\RaggedRight\arraybackslash}p{#1}}

\newcommand{\code}[1]{\nolinkurl{#1}}
\newcommand{\BA}{\textit{Base Analysis}}
\newcommand{\BAR}{\textit{BAReferent}}
\newcommand{\BAP}{\textit{BAProposition}}

\begin{document}

\begin{titlepage}
\centering
\vspace*{1.5cm}
{\Large Documentation-Driven Threat Analysis\par}
\vspace{0.8cm}
{\Huge\bfseries Documentation and Base Analysis Authoring Guide\par}
\vspace{0.25cm}
{\LARGE Revision 3 - governed documentation to BA handoff\par}
\vspace{0.9cm}
{\large Metodo operativo per documentare il minimum sufficient governed meaning e consegnarlo alla Base Analysis\par}
\vfill
\begin{minipage}{0.88\textwidth}
\small
\textbf{Stato.} Consolidamento human-readable successivo alla chiusura BA6 sul corpus governato Facial Access.

\medskip
\textbf{Project authority.} \code{FACIAL-ACCESS-GOV-R2}.

\medskip
\textbf{Accepted BA baseline.} \code{FACIAL-ACCESS-BA-R24-R1}.

\medskip
\textbf{Pinned source revision.} \code{8af2257a1df94fa5a83d4853ed0a1eb4d020c429}.

\medskip
Questa guida non sostituisce i contratti normativi BA0--BA6 né i documenti governati del progetto. Li organizza in una procedura leggibile e operativa.
\end{minipage}
\vfill
{\large 29 agosto 2026\par}
\end{titlepage}

\pagenumbering{roman}
\tableofcontents
\newpage
\pagenumbering{arabic}

\section{Scopo della guida}

Questa guida spiega come produrre documentazione governata che rimanga leggibile per le persone, stabile nel tempo e sufficientemente precisa per una successiva Base Analysis senza anticipare dettagli di implementazione non ancora decisi.

\begin{closedbox}[Principio centrale]
La documentazione deve governare il \textbf{minimum sufficient governed meaning}: abbastanza semantica per distinguere responsabilità, decisioni, comportamenti, condizioni e proprietà materialmente rilevanti, ma non più dettaglio di quello giustificato dal progetto.
\end{closedbox}

La completezza documentale non coincide con la descrizione esaustiva del sistema. Una baseline può essere adeguata anche quando alcuni aspetti restano \code{NOT SPECIFIED}, purché l'assenza non venga mascherata e non impedisca il significato necessario allo scope corrente.

\subsection{Il livello di dettaglio non è fissato globalmente}

DDTA non impone un livello uniforme come ``architettura'', ``componente'', ``interfaccia'' o ``protocollo''. Il criterio è adattivo:

\begin{lstlisting}
documentazione governata
    -> domanda semantica o di sicurezza
    -> l'evidenza e' sufficiente?
         YES -> stop
         NO  -> gap / NOT SPECIFIED / ulteriore review
\end{lstlisting}

La domanda di arresto è:

\begin{conceptbox}[Stopping criterion]
Abbiamo abbastanza semantica per formulare la proprietà o la domanda rilevante e dire, con evidenza documentale, se è soddisfatta, violata oppure non specificata?
\end{conceptbox}

\section{Authority prima della scrittura}

La cronologia non determina authority.

Prima di creare, revisionare o usare un documento:

\begin{enumerate}
\item identificare la baseline governata corrente;
\item distinguere artifact correnti, candidate, working e superseded;
\item preservare la provenance delle revisioni storiche;
\item evitare che BA, tool o analisi downstream diventino project authority.
\end{enumerate}

\begin{lstlisting}
project documentation
    = project authority

Base Analysis
    = analytical representation

threat method / test / projection
    = downstream consumer
\end{lstlisting}

Un finding downstream può proporre una clarification o una correzione, ma la project truth cambia soltanto attraverso review e aggiornamento della documentazione governata.

\section{Problem framing}

La documentazione inizia dalla responsabilità del progetto, non dalla tecnologia disponibile.

\begin{badbox}
\begin{lstlisting}
Il sistema usa una telecamera IP Ethernet,
un PC locale e un modello facciale per aprire il varco.
\end{lstlisting}
Questa frase incorpora già soluzione, placement, mezzo e attuazione fisica.
\end{badbox}

\begin{goodbox}
\begin{lstlisting}
Il progetto controlla un tentativo di accesso,
distingue la persona presente,
lo stato di autorizzazione rilevante
e la decisione sul singolo tentativo.
\end{lstlisting}
\end{goodbox}

Il framing è una disciplina di metodo. Non richiede necessariamente un nuovo tipo di documento.

\section{MacroRequirement}

Un MacroRequirement governa una responsabilità stabile del problema.

Deve rispondere almeno a:

\begin{itemize}
\item che cosa è responsabilità del branch;
\item quale significato o risultato macro produce/contribuisce;
\item quali responsabilità importanti sono fuori scope;
\item quali stakeholder sono materialmente interessati;
\item quali dipendenze macro esistono senza confondere dependency e containment.
\end{itemize}

Corpo consigliato:

\begin{lstlisting}
Title
Intent
Context
Stakeholders
Scope
Assumptions / Constraints   [solo se branch-wide]
dependsOn                   [solo quando genuino]
\end{lstlisting}

\subsection{Semantic-sufficiency gate}

Prima di scendere a Decision o Functional Requirement, verificare che l'MR abbia un solo significato stabile.

\begin{examplebox}[MR-0003]
La formulazione corrente chiarisce che, quando la responsabilità inizia, la persona è presente ma la specifica \code{GovernedIdentity} corrispondente non è ancora determinata.

Questa frase elimina la lettura alternativa in cui una identità specifica sarebbe già preselezionata e il sistema dovrebbe soltanto verificarla.
\end{examplebox}

\section{Decision}

Una Decision governa una scelta progettuale materialmente rilevante che specializza una responsabilità più stabile.

Una Decision dovrebbe esplicitare:

\begin{lstlisting}
Context
Decision
Consequences
\end{lstlisting}

Una scelta merita una Decision autonoma quando:

\begin{itemize}
\item può cambiare senza cambiare la responsabilità macro;
\item modifica placement, boundary, strategia o altro significato riusabile;
\item è necessario distinguerla da alternative progettuali reali;
\item il cambiamento può richiedere revalidazione downstream.
\end{itemize}

\begin{examplebox}[D-3.1, D-3.4, D-3.5, D-3.6]
Nel caso Facial Access sono Decision distinte:
\begin{itemize}
\item riconoscimento facciale come strategia;
\item separazione CameraSubsystem / RecognitionProcessor;
\item consumo di un servizio di connettività non posseduto dal progetto;
\item Ethernet cablata come mezzo corrente dell'interazione.
\end{itemize}
Separarle evita che ``facial recognition'', ``camera'', ``servizio'' ed ``Ethernet'' diventino un unico commitment indivisibile.
\end{examplebox}

\section{Functional Requirement}

Un Functional Requirement governa un comportamento operativo verificabile derivato da una responsabilità/Decision.

Una buona formulazione identifica:

\begin{itemize}
\item quando il comportamento è richiesto;
\item chi ne possiede la responsabilità funzionale;
\item il significato di input/output o stato coinvolto;
\item il binding contestuale necessario;
\item ciò che il FR deliberatamente non governa.
\end{itemize}

\subsection{Non anticipare la BA}

La documentazione deve governare il significato del progetto, non scegliere un operatore BA per comodità.

\begin{examplebox}[FR-3.4.1 - acquire]
La fonte governa che \code{CameraSubsystem} acquisisca una nuova \code{RecognitionCapture} quando la \code{IdentityDeterminationRequest} la richiede.

La fonte non decide se ``acquire'' debba essere modellato come:
\begin{lstlisting}
create
observe
produce
combination
\end{lstlisting}
Questa under-specification è intenzionale e corretta.
\end{examplebox}

\subsection{Binding contestuale}

Quando due significati devono rimanere associati alla stessa richiesta, identità, tentativo o valutazione, il binding appartiene alla semantica del FR se è necessario per evitare un risultato materialmente diverso.

Nel caso Facial Access:

\begin{lstlisting}
RecognitionCapture
    <-> IdentityDeterminationRequest
\end{lstlisting}

è parte dell'obbligo di acquisizione e deve essere preservato nel delivery.

\section{Requirement split e unità coerente}

Non dividere automaticamente un requisito in una riga per verbo. Dividere quando esistono obblighi con:

\begin{itemize}
\item semantic owner diverso;
\item failure mode materialmente distinto;
\item lifecycle/review indipendente;
\item diversa specializzazione di sicurezza;
\item diversa possibilità di change/revalidation.
\end{itemize}

Evitare anche l'estremo opposto: un FR che contiene strategia, algoritmo, trasporto, security mechanism e policy organizzativa in un unico testo.

\section{Specialized Requirement}

Uno Specialized Requirement è giustificato quando una proprietà più specifica deve essere governata autonomamente rispetto al FR parent.

Non deve essere creato solo per riempire una gerarchia.

\begin{examplebox}[DG-FA-002]
La review Facial Access non ha trovato evidenza sufficiente per introdurre un requisito quantitativo di ``capture quality''.

Il significato necessario è stato governato nell'esito:
\begin{lstlisting}
SUCCESS
NEGATIVE
INCONCLUSIVE
\end{lstlisting}

Non sono stati inventati score, confidence, threshold o ranking.
\end{examplebox}

\section{Security Requirement}

Un Security Requirement deve specializzare un comportamento o significato governato e mantenere esplicito lo scope della proprietà.

Nel corpus Facial Access:

\begin{lstlisting}
FR-3.4.2 RecognitionCaptureDelivery
    -> Confidentiality
    -> Integrity
    -> AuthorizedProvenance
\end{lstlisting}

Queste proprietà non vengono propagate automaticamente a tutto il sistema.

\subsection{Proprietà non significa meccanismo}

\begin{badbox}
\begin{lstlisting}
Confidentiality -> TLS
Integrity       -> signature
Provenance      -> certificate
\end{lstlisting}
Nessuna di queste equivalenze è autorizzata se il progetto non governa il meccanismo.
\end{badbox}

\subsection{NOT SPECIFIED può essere corretto}

Il corpus governa le proprietà C/I/AuthorizedProvenance ma non la base completa di autorizzazione per:

\begin{itemize}
\item chi possa conoscere il contenuto della capture (\code{AUTH-C});
\item quale evidenza renda una origine autorizzata per una specifica request (\code{AUTH-P}).
\end{itemize}

Questi significati restano \code{NOT SPECIFIED} e non bloccano la baseline dimostrativa corrente.

\section{Cross-MR semantics}

Quando un risultato attraversa responsabilità diverse, il downstream consumer deve essere governato per responsabilità, non inventato come componente tecnico.

Nel caso corrente:

\begin{lstlisting}
MR-0003
    -> IdentityDeterminationOutcome
       -> SUCCESS makes GovernedIdentity X available

MR-0002
    -> AccessAuthorizationState concerning GovernedIdentity X

MR-0001
    -> uses both meanings
       -> AccessDecision
\end{lstlisting}

L'identità usata per l'autorizzazione deve essere la stessa resa disponibile dalla determinazione riuscita.

\section{Consumed service boundary}

Un servizio consumato non diventa automaticamente responsabilità del progetto.

\begin{examplebox}[D-3.5]
La realizzazione corrente consuma \code{ConnectivityService}. Il progetto non possiede né gestisce l'infrastruttura sottostante.

Quindi:
\begin{lstlisting}
consumption != ownership
requirement  != provider guarantee
wired medium != hidden end-to-end topology
\end{lstlisting}
\end{examplebox}

Se una proprietà richiesta al servizio non è garantita da evidenza governata, il risultato è una assurance/coverage question, non una nuova project truth inventata.

\section{Decisioni con condizioni congiunte}

Una policy decisionale deve governare abbastanza semantica da distinguere condizioni necessarie, ramo positivo e rami non consentiti.

\begin{examplebox}[D-1.1 / FR-1.1]
\begin{lstlisting}
SUCCESS
AND required authorization condition satisfied
    for the same GovernedIdentity
    -> AccessDecision MUST ALLOW

NEGATIVE
    -> MUST NOT ALLOW

INCONCLUSIVE
    -> MUST NOT ALLOW

SUCCESS
AND required authorization condition not satisfied
    -> MUST NOT ALLOW
\end{lstlisting}
\end{examplebox}

La documentazione non è obbligata a introdurre un campo \code{authorized = TRUE}. Se quel vocabolario non esiste, non deve essere creato per facilitare la BA.

\section{Documentation gaps e clarification}

Un gap è una diagnosi, non un requisito automatico.

Classificazioni utili:

\begin{lstlisting}
missing governed meaning
conflicting governed meaning
ambiguous semantic owner
consumer evidence insufficient
analysis-relevant NOT SPECIFIED
deferred outside current scope
\end{lstlisting}

Il ciclo corretto è:

\begin{lstlisting}
analysis / BA question
    -> diagnostic
    -> semantic-owner review
    -> project-document change only if justified
    -> rebuild BA
\end{lstlisting}

\section{Canonical terminology}

I riferimenti semanticamente operativi devono avere identità/nome canonico controllato.

Una frase umana può dire ``la camera invia la capture'', ma questo non crea automaticamente alias normativi come:

\begin{lstlisting}
send
deliver
push
tx
\end{lstlisting}

per un eventuale operatore BA.

Lo stesso vale per nomi di referent e valori controllati.

\section{Promotion gate}

Una baseline documentale è pronta per promotion quando:

\begin{itemize}
\item il semantic owner dei significati materiali è stabilito;
\item le cross-responsibility dependencies necessarie sono governate;
\item i gap bloccanti sono risolti o esplicitamente esclusi;
\item i \code{NOT SPECIFIED} non bloccanti sono preservati;
\item una regression semantica completa non rileva contraddizioni non gestite;
\item la governance autorizza esplicitamente l'uso come source per BA.
\end{itemize}

\section{Handoff alla Base Analysis}

Il passaggio alla BA non significa ``tradurre ogni frase''.

L'handoff fornisce:

\begin{itemize}
\item baseline authority key;
\item pinned source revision;
\item source set;
\item termini/referent governati;
\item relazioni/condizioni materialmente necessarie;
\item diagnostics e \code{NOT SPECIFIED};
\item change/revalidation context quando rilevante.
\end{itemize}

La BA applica poi il proprio criterio minimo e i contratti BA0--BA6.

\section{Facial Access - esempio end-to-end}

La documentazione corrente dell'esempio è:

\begin{lstlisting}
FACIAL-ACCESS-GOV-R2
CURRENT_GOVERNED

MR-0001 Access Control
MR-0002 Access Authorization
MR-0003 Identity Determination
\end{lstlisting}

I passaggi principali sono:

\begin{enumerate}
\item MR-0003 chiarisce che l'identità specifica non è preselezionata.
\item D-3.1 sceglie riconoscimento facciale.
\item D-3.2 governa SUCCESS/NEGATIVE/INCONCLUSIVE.
\item D-3.4 separa acquisizione e riconoscimento.
\item FR-3.4.1 governa acquire senza specificare snapshot/streaming.
\item FR-3.4.2 governa delivery e request binding.
\item D-3.5 governa servizio consumato.
\item D-3.6 governa Ethernet cablata sul delivery corrente.
\item SEC-3.4.2-C/I/P specializzano quel delivery.
\item MR-0002 governa stato di autorizzazione riferito a GovernedIdentity senza enum interno.
\item D-1.1/FR-1.1 governano la policy congiuntiva e il ramo positivo.
\item la semantic regression chiude il corpus.
\item la governance promuove \code{FACIAL-ACCESS-GOV-R2}.
\item la BA viene costruita e regressata separatamente.
\end{enumerate}

\section{Checklist pratica di authoring}

Prima di chiudere una revisione chiedere:

\begin{enumerate}
\item Sto governando un problema/responsabilità o una soluzione prematura?
\item Il semantic owner è unico e chiaro?
\item Ho separato Decision e comportamento funzionale?
\item Ho introdotto dettagli che il progetto non ha deciso?
\item Un binding di request/identity/context è materialmente necessario?
\item Una proprietà è scoped al comportamento corretto?
\item Sto confondendo proprietà di sicurezza e meccanismo?
\item Sto trasformando un gap in requisito senza review?
\item Un \code{NOT SPECIFIED} può rimanere esplicito?
\item Un downstream consumer potrebbe ottenere due letture materialmente diverse?
\item Il cambiamento di una Decision può essere revalidato senza riscrivere l'MR?
\item La baseline è davvero autorizzata per handoff BA?
\end{enumerate}

\section{Cosa non deve fare la documentazione}

La documentazione DDTA non deve:

\begin{itemize}
\item anticipare STRIDE/STRIDE-AI;
\item creare un DFD solo perché un threat tool lo richiede;
\item descrivere provider hop non governati;
\item introdurre protocolli e security mechanisms per default;
\item scegliere operatori BA;
\item creare sinonimi normativi non registrati;
\item nascondere ambiguità con wording generico;
\item trasformare automaticamente findings in project truth.
\end{itemize}

\section{Stato corrente e prossimo uso}

La baseline governata dell'esempio è \code{FACIAL-ACCESS-GOV-R2}. La Base Analysis accettata per lo scope R24 è \code{FACIAL-ACCESS-BA-R24-R1}.

Questa guida R3 costituisce la guida human-readable corrente per la fase documentation-to-BA.

Il prossimo uso previsto è la revisione dei capitoli di tesi che descrivono:

\begin{itemize}
\item progressive governed documentation;
\item semantic sufficiency;
\item minimum sufficient governed meaning;
\item authority-preserving feedback;
\item handoff alla Base Analysis.
\end{itemize}

La valutazione STRIDE/STRIDE-AI resta una fase successiva.
\end{document}
